\documentclass[oneside,a4paper,14pt]{extarticle}
\usepackage[T1,T2A]{fontenc}
\usepackage[a4paper,letterpaper,top=20mm,bottom=20mm,left=30mm,right=10mm]{geometry}
\usepackage[utf8]{inputenc}
\usepackage[russian]{babel}
\usepackage{textcomp}
\usepackage{indentfirst}
\usepackage{graphicx}
\usepackage{mwe}
\usepackage{wrapfig}
\usepackage{caption}
\usepackage{amsmath}
\usepackage{amsfonts}
\usepackage{amsthm}
\usepackage[all]{xy}
\usepackage[breaklinks]{hyperref}
\usepackage{titlesec}
\titleformat{\section}{\normalsize\bfseries}{\thesection}{1em}{}
\titleformat{\subsection}{\normalsize\bfseries}{\thesubsection}{1em}{}
\titleformat{\subsubsection}{\normalsize\bfseries}{\thesubsection}{1em}{}
\renewcommand\baselinestretch{1.33}\normalsize
\setlength{\parindent}{1.25cm}
\usepackage{indentfirst}

\begin{document}
    \newpage\thispagestyle{empty}
    \begin{center}
        МИНИСТЕРСТВО НАУКИ И ВЫСШЕГО ОБРАЗОВАНИЯ\\
            РОССИЙСКОЙ ФЕДЕРАЦИИ
            ФЕДЕРАЛЬНОЕ ГОСУДАРСТВЕННОЕ БЮДЖЕТНОЕ\\
            ОБРАЗОВАТЕЛЬНОЕ
            УЧРЕЖДЕНИЕ ВЫСШЕГО ОБРАЗОВАНИЯ\\
            «ВЯТСКИЙ ГОСУДАРСТВЕННЫЙ УНИВЕРСИТЕТ»\\
            Институт математики и информационных систем\\
            Факультет автоматики и вычислительной техники\\
            Кафедра электронных вычислительных машин
    \end{center}
    \vspace{20mm}

    \begin{center}
        Отчёт по лабораторной работе №7\\
        по дисциплине\\
        <<Программирование>>\\
    \end{center}
    \vspace{48mm}

    Выполнил студент гр. ИВТб-1302-06-00 \hspace{10mm} \rule[-0,5mm]{30mm}{0.15mm}\,/Изместьев М.Н./


    Проверил доцент кафедры ЭВМ \hfill  \rule[-0,5mm]{30mm}{0.15mm}\,/Баташев П.А./

    \vfill
    \begin{center}
        Киров\\
        2026
    \end{center}

    \newpage\thispagestyle{empty}

\begin{flushleft}
\section*{Цель работы:}
\noindent Научиться работать с табличными данными при помощи библиотек \texttt{pandas} и \texttt{numpy}, строить статистические и аналитические графики с помощью \texttt{seaborn} и \texttt{matplotlib}, а также встраивать интерактивные графики в оконное приложение на \texttt{tkinter}.

\section*{Задание.}

1. Разработать интерактивный дашборд с использованием стека pandas + seaborn + matplotlib + tkinter.

2. Реализовать не менее четырёх типов графиков: линейный, столбчатый, точечный, тепловая карта.

3. Выполнить предобработку данных: фильтрацию физически невозможных значений, удаление выбросов (IQR), создание производных признаков, биннинг, скользящее среднее.

4. Реализовать панель фильтров: выбор спортсмена, зоны пульса, максимального темпа, окна сглаживания.

5. Обеспечить интерактивность: переключение типа графика, смена режима агрегации, экспорт в PNG/PDF.

6. В качестве ответа прикрепить отчёт и ссылку на репозиторий с проектом.

\textbf{Вариант 10} --- Спортивная аналитика.\\
\textbf{Стек:} Python, pandas, seaborn, matplotlib, tkinter, numpy.\\
\textbf{Данные:} \texttt{data.csv} --- треки спортивных тренировок. Столбцы: \texttt{athlete\_id}, \texttt{ts} (timestamp), \texttt{dist} (дистанция, км), \texttt{pace} (темп, мин/км), \texttt{cal} (калории), \texttt{zone} (пульсовая зона 1--5).

\section*{Описание программы}

Программа реализована в \textbf{процедурном стиле} (без классов, кроме \texttt{RussianToolbar}). Точка входа --- функция \texttt{main()}: загружает данные, строит GUI, запускает цикл событий \texttt{tkinter}.

\textbf{Этапы обработки данных:}

\begin{enumerate}
\item Загрузка CSV через \texttt{pd.read\_csv}. Поиск файла --- сначала в директории скрипта, затем в корне проекта.
\item Фильтрация: удаление строк с \texttt{NaN}/\texttt{Inf} в ключевых столбцах, отсечение отрицательных значений дистанции и темпа, замена отрицательных калорий на 0.
\item Динамические фильтры из UI: по \texttt{athlete\_id}, \texttt{zone}, максимальному темпу.
\item Производные признаки: \texttt{speed\_kmh = dist / (pace/60)}, \texttt{cal\_per\_km = cal / dist}.
\item Удаление выбросов: IQR-обрезка по группам (\texttt{athlete\_id}) через \texttt{groupby + transform}.
\item Категориальные типы и биннинг дистанции (\texttt{pd.cut}, 5 бинов).
\item Временная ось (\texttt{pd.to\_datetime}) и скользящее среднее (\texttt{df.rolling(k)}).
\end{enumerate}

\textbf{Типы графиков:}
\begin{itemize}
\item \textbf{Линейный} --- скользящая средняя скорости, агрегация по часам (\texttt{df.resample("h")}).
\item \textbf{Столбчатый} --- калории по зонам пульса (\texttt{groupby + mean/sum/median}).
\item \textbf{Точечный} --- темп vs калории/км, цвет --- пульсовая зона (\texttt{sns.scatterplot}).
\item \textbf{Тепловая карта} --- средняя скорость: зона × бин дистанции (\texttt{pd.pivot\_table + sns.heatmap}).
\end{itemize}
\end{flushleft}

\newpage
\textbf{Загрузка данных и предобработка (dashboard.py)}
\small
\begin{verbatim}
import tkinter as tk
from tkinter import filedialog, messagebox, ttk
from pathlib import Path
import matplotlib.pyplot as plt
import numpy as np
import pandas as pd
import seaborn as sns
from matplotlib.backends.backend_tkagg import (
    FigureCanvasTkAgg, NavigationToolbar2Tk)

df_raw  = None
df_work = None
fig     = plt.Figure(figsize=(10, 6), dpi=100)
canvas  = None
root    = None
current_chart = "line"
agg_mode      = "mean"
AGG_LABELS = {"mean": "среднее", "sum": "сумма", "median": "медиана"}


def load_dataset() -> pd.DataFrame:
    local  = Path("data.csv")
    parent = Path(__file__).resolve().parent.parent / "data.csv"
    source = local if local.exists() else parent
    if not source.exists():
        raise FileNotFoundError("Файл data.csv не найден")
    return pd.read_csv(source)


def grouped_iqr_clip(series: pd.Series,
                     group_key: pd.Series) -> pd.Series:
    tmp = pd.DataFrame({"x": series, "g": group_key})
    q1  = tmp.groupby("g")["x"].transform(lambda s: s.quantile(0.25))
    q3  = tmp.groupby("g")["x"].transform(lambda s: s.quantile(0.75))
    iqr = q3 - q1
    return tmp["x"].clip(lower=q1 - 1.5 * iqr, upper=q3 + 1.5 * iqr)


def preprocess_data() -> pd.DataFrame:
    global df_work
    work = df_raw.copy()

    # 1. Фильтрация физически невозможных значений
    work = work[
        np.isfinite(work["dist"]) &
        np.isfinite(work["pace"]) &
        np.isfinite(work["cal"])
    ].copy()
    work = work[(work["dist"] >= 0) & (work["pace"] >= 0)].copy()
    work["cal"] = np.where(work["cal"] < 0, 0.0, work["cal"])

    # 2. Динамические фильтры из UI
    selected_athlete = athlete_var.get()
    selected_zone    = zone_var.get()
    max_pace         = float(pace_limit_var.get())
    if selected_athlete != "Все":
        work = work.loc[work["athlete_id"]==int(selected_athlete)].copy()
    if selected_zone != "Все":
        work = work.loc[work["zone"]==int(selected_zone)].copy()
    work = work.loc[work["pace"] <= max_pace].copy()

    # 3. Производные признаки
    work["speed_kmh"]  = work["dist"] / np.maximum(work["pace"]/60.0, 1e-8)
    work["cal_per_km"] = work["cal"]  / np.maximum(work["dist"], 1e-8)

    # 4. IQR-обрезка по группам (athlete_id)
    work["pace"]       = grouped_iqr_clip(work["pace"],work["athlete_id"])
    work["cal_per_km"] = grouped_iqr_clip(work["cal_per_km"],
                                          work["athlete_id"])

    # 5. Категории + биннинг
    work["athlete_id"] = work["athlete_id"].astype("category")
    work["zone"]       = work["zone"].astype("category")
    work["dist_bin"]   = pd.cut(
        work["dist"],
        bins=[0, 2, 5, 10, 20, float("inf")],
        labels=["<2 км","2-5 км","5-10 км","10-20 км",">20 км"],
        right=False)

    # 6. Временная ось + скользящее среднее
    work["dt"] = pd.to_datetime(work["ts"], unit="s", errors="coerce")
    work = work.dropna(subset=["dt"]).sort_values("dt")
    k = max(2, int(roll_window_var.get()))
    work["speed_roll"] = work["speed_kmh"].rolling(k, min_periods=1).mean()

    df_work = work
    return work
\end{verbatim}

\newpage
\textbf{Функции отрисовки графиков}
\small
\begin{verbatim}
def draw_line() -> None:
    fig.clf()
    ax   = fig.add_subplot(111)
    data = (df_work.set_index("dt")
                   .resample("h")["speed_roll"].mean().dropna())
    sns.lineplot(x=data.index, y=data.values, ax=ax, color="#2E86DE")
    ax.set_title("Скользящая средняя скорости (почасово)")
    ax.set_xlabel("Время"); ax.set_ylabel("Скорость, км/ч")
    ax.grid(True, alpha=0.2)
    fig.tight_layout(); canvas.draw_idle()


def draw_bar() -> None:
    fig.clf()
    ax      = fig.add_subplot(111)
    grouped = df_work.groupby("zone", observed=True)["cal"]
    if   agg_mode == "sum":    s = grouped.sum()
    elif agg_mode == "median": s = grouped.median()
    else:                      s = grouped.mean()
    sns.barplot(x=s.index.astype(str), y=s.values,
                ax=ax, palette="viridis", errorbar=None)
    ax.set_title(f"Калории по зонам пульса ({AGG_LABELS[agg_mode]})")
    ax.set_xlabel("Зона пульса"); ax.set_ylabel("Калории")
    fig.tight_layout(); canvas.draw_idle()


def draw_scatter() -> None:
    fig.clf()
    ax     = fig.add_subplot(111)
    sample = df_work.sample(min(2500, len(df_work)), random_state=42)
    sns.scatterplot(data=sample, x="pace", y="cal_per_km",
                    hue="zone", alpha=0.55, s=25, ax=ax)
    ax.set_title("Темп и калории на километр")
    ax.set_xlabel("Темп, мин/км"); ax.set_ylabel("Калории/км")
    fig.tight_layout(); canvas.draw_idle()


def draw_heatmap() -> None:
    fig.clf()
    ax    = fig.add_subplot(111)
    pivot = pd.pivot_table(df_work, index="zone", columns="dist_bin",
                           values="speed_kmh", aggfunc="mean", observed=True)
    if pivot.empty:
        ax.text(0.5, 0.5, "Нет данных", ha="center", va="center")
        ax.axis("off")
    else:
        sns.heatmap(pivot, cmap="coolwarm", ax=ax,
                    cbar_kws={"label": "км/ч"}, annot=False)
    ax.set_title("Средняя скорость: зона x дистанция")
    ax.set_xlabel("Диапазон дистанции"); ax.set_ylabel("Пульсовая зона")
    fig.tight_layout(); canvas.draw_idle()
\end{verbatim}

\newpage
\textbf{Построение GUI (build\_ui)}
\small
\begin{verbatim}
def build_ui() -> None:
    global root, canvas, athlete_var, zone_var, \
           pace_limit_var, roll_window_var, status_var

    sns.set_theme(style="whitegrid")
    plt.rcParams["font.family"]        = ["Segoe UI", "Arial", "DejaVu Sans"]
    plt.rcParams["axes.unicode_minus"] = False

    root = tk.Tk()
    root.title("Лаб. 7 -- Интерактивный дашборд (Вариант 10)")
    root.geometry("1280x820")
    root.configure(bg="#f3f5f7")

    athlete_var     = tk.StringVar(value="Все")
    zone_var        = tk.StringVar(value="Все")
    pace_limit_var  = tk.DoubleVar(value=12.0)
    roll_window_var = tk.IntVar(value=10)
    status_var      = tk.StringVar(value="Готово")

    athlete_values = ["Все"] + sorted(
        df_raw["athlete_id"].dropna().astype(int).astype(str).unique())
    zone_values    = ["Все"] + sorted(
        df_raw["zone"].dropna().astype(int).astype(str).unique())

    top = tk.Frame(root, bg="#f3f5f7")
    top.pack(side=tk.TOP, fill=tk.X, padx=12, pady=8)
    tk.Label(top, text="Спортсмен:", bg="#f3f5f7").pack(side=tk.LEFT)
    ttk.Combobox(top, values=athlete_values, textvariable=athlete_var,
                 width=8, state="readonly").pack(side=tk.LEFT, padx=6)
    tk.Label(top, text="Зона:", bg="#f3f5f7").pack(side=tk.LEFT)
    ttk.Combobox(top, values=zone_values, textvariable=zone_var,
                 width=6, state="readonly").pack(side=tk.LEFT, padx=6)
    tk.Label(top, text="Макс. темп:", bg="#f3f5f7").pack(side=tk.LEFT, padx=(12,2))
    tk.Scale(top, variable=pace_limit_var, from_=3.0, to=15.0,
             resolution=0.1, orient=tk.HORIZONTAL, length=190,
             bg="#f3f5f7").pack(side=tk.LEFT)
    tk.Label(top, text="Окно k:", bg="#f3f5f7").pack(side=tk.LEFT, padx=(12,2))
    tk.Scale(top, variable=roll_window_var, from_=2, to=60,
             resolution=1, orient=tk.HORIZONTAL, length=160,
             bg="#f3f5f7").pack(side=tk.LEFT)
    tk.Button(top, text="  Экспорт  ",
              command=export_plot).pack(side=tk.RIGHT, padx=4)
    tk.Button(top, text="  Обновить  ", command=refresh_data,
              bg="#2E86DE", fg="white", relief=tk.FLAT).pack(
              side=tk.RIGHT, padx=4)

    mid = tk.Frame(root, bg="#e8edf2", pady=5)
    mid.pack(side=tk.TOP, fill=tk.X)
    tk.Label(mid, text="График:", bg="#e8edf2",
             font=("Segoe UI",9,"bold")).pack(side=tk.LEFT, padx=(12,4))
    for label, key in [("Линейный","line"),("Столбчатый","bar"),
                        ("Точечный","scatter"),("Тепловая карта","heatmap")]:
        tk.Button(mid, text=label, width=14,
                  command=lambda k=key: set_chart(k)).pack(side=tk.LEFT, padx=3)
    tk.Label(mid, text="  Агрегация:", bg="#e8edf2",
             font=("Segoe UI",9,"bold")).pack(side=tk.LEFT, padx=(16,4))
    for label, val in [("Среднее","mean"),("Сумма","sum"),("Медиана","median")]:
        tk.Radiobutton(mid, text=label, value=val,
                       command=lambda v=val: set_agg(v),
                       bg="#e8edf2").pack(side=tk.LEFT)

    plot_frame = tk.Frame(root, bg="white", relief=tk.SUNKEN, bd=1)
    plot_frame.pack(fill=tk.BOTH, expand=True, padx=12, pady=6)
    canvas_widget = FigureCanvasTkAgg(fig, master=plot_frame)
    canvas_widget.get_tk_widget().pack(fill=tk.BOTH, expand=True)
    global canvas
    canvas = canvas_widget
    toolbar = RussianToolbar(canvas, plot_frame)
    toolbar.update()
    toolbar.pack(side=tk.TOP, fill=tk.X)

    tk.Label(root, textvariable=status_var, anchor="w",
             bg="#d0d7df", font=("Segoe UI",9)).pack(
        side=tk.BOTTOM, fill=tk.X, ipady=3)


def main() -> None:
    global df_raw
    df_raw = load_dataset()
    build_ui()
    refresh_data()
    root.mainloop()

if __name__ == "__main__":
    main()
\end{verbatim}

\normalsize
\newpage
\begin{flushleft}
\section*{Используемые методы pandas и seaborn}

\begin{center}
\begin{tabular}{|l|l|}
\hline
\textbf{Метод} & \textbf{Применение} \\
\hline
\texttt{pd.read\_csv} & Загрузка датасета \\
\texttt{df.copy()} & Изоляция изменений от исходника \\
\texttt{np.isfinite} & Фильтрация NaN и Inf \\
\texttt{np.where} & Векторная замена отрицательных калорий \\
\texttt{df.loc[mask]} & Динамические фильтры из UI \\
\texttt{np.maximum} & Защита от деления на ноль \\
\texttt{groupby + transform} & IQR-обрезка по группам \\
\texttt{df.clip} & Ограничение выбросов \\
\texttt{astype("category")} & Оптимизация памяти \\
\texttt{pd.cut} & Биннинг дистанции на 5 интервалов \\
\texttt{pd.to\_datetime} & Парсинг UNIX-timestamp \\
\texttt{df.rolling(k).mean()} & Скользящее среднее скорости \\
\texttt{df.resample("h")} & Почасовая агрегация для линейного графика \\
\texttt{df.groupby + agg} & Калории по зонам для столбчатого \\
\texttt{df.sample} & Случайная выборка для scatter (2500 точек) \\
\texttt{pd.pivot\_table} & Сводная таблица для тепловой карты \\
\texttt{sns.lineplot} & Линейный график с seaborn \\
\texttt{sns.barplot} & Столбчатый график \\
\texttt{sns.scatterplot} & Точечный с цветом по зоне \\
\texttt{sns.heatmap} & Тепловая карта \\
\hline
\end{tabular}
\end{center}
\end{flushleft}

\newpage
\textbf{Экранные формы:}

\centering
    \includegraphics[width=0.85\linewidth]{1.png}
    \label{fig:screen1}

    \textbf{Рисунок 1 --- }Главное окно дашборда, линейный график скользящей средней скорости

\centering
    \includegraphics[width=0.85\linewidth]{2.png}
    \label{fig:screen2}

    \textbf{Рисунок 2 --- }Столбчатый график калорий по зонам пульса (режим «Среднее»)

\centering
    \includegraphics[width=0.85\linewidth]{3.png}
    \label{fig:screen3}

    \textbf{Рисунок 3 --- }Точечный график «Темп vs Калории/км» с цветовой кодировкой зон

\centering
    \includegraphics[width=0.85\linewidth]{4.png}
    \label{fig:screen4}

    \textbf{Рисунок 4 --- }Тепловая карта средней скорости: пульсовая зона × диапазон дистанции

\centering
    \includegraphics[width=0.85\linewidth]{5.png}
    \label{fig:screen5}

    \textbf{Рисунок 5 --- }Панель фильтров: фильтрация по спортсмену и зоне пульса

\begin{flushleft}
\newpage
\section*{Вывод}
\noindent
В ходе лабораторной работы был создан интерактивный дашборд на стеке pandas + seaborn + matplotlib + tkinter. Реализованы полный цикл предобработки данных (фильтрация, IQR-обрезка, feature engineering, биннинг, rolling mean) и четыре типа графиков (линейный, столбчатый, точечный, тепловая карта). Панель фильтров позволяет динамически менять параметры визуализации без перезапуска программы. Кастомный тулбар \texttt{RussianToolbar} локализован на русский язык. Экспорт графика доступен в форматах PNG и PDF.
\end{flushleft}

\end{document}
